\documentclass[a4paper]{article}

\usepackage[english]{babel}
\usepackage[T1]{fontenc}

\usepackage[a4paper,top=3cm,bottom=2cm,left=3cm,right=3cm,marginparwidth=1.75cm]{geometry}

\usepackage{amsmath}
\usepackage{graphicx}
\usepackage[colorinlistoftodos]{todonotes}
\usepackage[colorlinks=true, allcolors=blue]{hyperref}

\title{Your Paper}
\author{You}

\begin{document}

\maketitle

\textbf{1. (A Simple Delegated Monitor Model) Consider a firm who needs
\$100,000 to initiate a project. The payout of this project is \$200,000 with probability 90\%, \$100,000 with probability 5\% and \$0 with probability 5\% at the end of two years. Suppose that there are 1000 small
investors each with \$100 to lend. If the project is
initiated, only the firm manager knows the result of the
project, while the lenders cannot see the results. However, as long as one of the lenders pay monitor fee of \$200, the results of the project will be known to everyone.}

\hypertarget{a}{%
\subsection{1a)}\label{a}}

Assume each investor knows that there are many other small investors
with \$100 to lend, but each investor lives in an isolated island and
cannot coordinate with other investors, he or she can only communicate
with the firm. Suppose the project is the only investment opportunity
for the small investors. Will the small investors lend directly to the
firm? Will the small investor pay the monitor fee? Why or Why not?
(Hint: Solve this question backward in two steps. First, suppose
everyone invests \$100, will they pay monitoring cost? Draw the (gross)
payout table and solve the Nash Equilibrium. Then, find out whether the
smaller investors want to invest or not by comparing the payout of
investing and not investing.

\textbf{Step 1: Will investors pay the monitoring fee? (Given everyone
invested)}

Let\textquotesingle s construct the Nash Equilibrium payout table for
the monitoring decision:

\textbf{Payout Matrix for One Investor:}

\begin{longtable}[]{@{}
  >{\raggedright\arraybackslash}p{(\columnwidth - 4\tabcolsep) * \real{0.2196}}
  >{\raggedright\arraybackslash}p{(\columnwidth - 4\tabcolsep) * \real{0.3890}}
  >{\raggedright\arraybackslash}p{(\columnwidth - 4\tabcolsep) * \real{0.3913}}@{}}
\toprule()
\begin{minipage}[b]{\linewidth}\raggedright
\textbf{Your Decision}
\end{minipage} & \begin{minipage}[b]{\linewidth}\raggedright
\textbf{Others Monitor}
\end{minipage} & \begin{minipage}[b]{\linewidth}\raggedright
\textbf{Others Don\textquotesingle t Monitor}
\end{minipage} \\
\midrule()
\endhead
\textbf{Monitor} & Your share of payout - \$200 & Your share of payout -
\$200 \\
\textbf{Don\textquotesingle t Monitor} & Your share of payout - \$0 &
Firm claims \$0, you get \$0 \\
\bottomrule()
\end{longtable}

\textbf{Analysis:}

\begin{itemize}
\item
  \textbf{If others monitor}: Paying \$200 is unnecessary since results
  are already revealed (dominant strategy: Don\textquotesingle t
  Monitor)
\item
  \textbf{If others don\textquotesingle t monitor}: Without monitoring,
  the firm manager can claim the project failed (payout = \$0) and keep
  all profits. Since monitoring costs \$200 but your expected share of
  the project is only:

  \begin{itemize}
  \item
    Investment per investor: \$100
  \item
    Your share = \$100/\$100,000 = 0.1\% of total payout
  \item
    Expected share = 0.001 × (\$200,000 × 0.9 + \$100,000 × 0.05 + \$0 ×
    0.05) = 0.001 × \$185,000 = \$185
  \end{itemize}
\end{itemize}

But if you\textquotesingle re the only one monitoring, you pay \$200 to
reveal the truth, but the firm manager still has no obligation to pay
fairly without enforcement. Even worse, the free-rider problem means no
single investor wants to bear the full \$200 cost.

\textbf{Nash Equilibrium: Nobody monitors}

\textbf{Step 2: Will investors lend to the firm?}

Given that nobody will monitor (from Step 1):

\begin{itemize}
\item
  Without monitoring, the firm manager always claims the project failed
  (payout = \$0)
\item
  Rational firm manager keeps all proceeds
\item
  Expected return for investor = \$0
\item
  Cost of lending = \$100
\end{itemize}

\textbf{Comparison:}

\begin{itemize}
\item
  Return from lending: \$0 (with certainty, due to information
  asymmetry)
\item
  Return from not lending: \$100 (keep the money)
\end{itemize}

\textbf{Answer: NO, small investors will NOT lend to the firm}

\textbf{Why?}

\begin{enumerate}
\def\labelenumi{\arabic{enumi}.}
\item
  Classic free-rider problem: Each investor hopes others will pay the
  \$200 monitoring fee
\item
  Monitoring cost (\$200) exceeds individual benefit (\$185 expected
  share)
\item
  Without monitoring, the firm manager has no incentive to report
  honestly
\item
  Rational expectation: zero return on investment
\item
  Outcome: Market failure due to information asymmetry and coordination
  failure
\end{enumerate}

\hypertarget{b}{%
\subsection{1b)}\label{b}}

\textbf{Will the project be funded if 1000 small investors can
communicate with each other? If so, explain how can they do it and what
is expected rate of return on the investment?}

\textbf{YES, the project will be funded}

\textbf{How they can structure it:}

\textbf{Option 1: Shared Monitoring Cost}

\begin{itemize}
\item
  All 1,000 investors agree to contribute \$0.20 each for monitoring
  (\$200 total)
\item
  Each investor\textquotesingle s total cost: \$100 (loan) + \$0.20
  (monitoring) = \$100.20
\item
  With monitoring, they can verify results and enforce repayment
\end{itemize}

\textbf{Option 2: Delegated Monitor (Financial Intermediary)}

\begin{itemize}
\item
  Investors pool funds and create a bank/intermediary
\item
  The intermediary collects \$100,000, pays \$200 monitoring fee
\item
  Acts as a single lender with monitoring capability
\item
  This demonstrates the \textbf{fundamental role of banks as delegated
  monitors}
\end{itemize}

\textbf{Expected Rate of Return Calculation:}

\textbf{Gross Expected Payout:}

\begin{itemize}
\item
  E{[}Payout{]} = \$200,000 × 0.9 + \$100,000 × 0.05 + \$0 × 0.05 =
  \$185,000
\end{itemize}

\textbf{Total Investment:}

\begin{itemize}
\item
  Principal: \$100,000
\item
  Monitoring cost: \$200
\item
  Total cost: \$100,200
\end{itemize}

\textbf{Net Profit:}

\begin{itemize}
\item
  Net = \$185,000 - \$100,200 = \$84,800
\end{itemize}

\textbf{Expected Return over 2 years:}

\begin{itemize}
\item
  Return = \$84,800 / \$100,200 = 84.63\% over 2 years
\item
  Annualized return ≈ 36.0\% per year (using compound growth:
  (1.8463)\^{}0.5 - 1)
\end{itemize}

\textbf{Per Investor:}

\begin{itemize}
\item
  Each invested: \$100.20
\item
  Expected return after 2 years: \$100.20 × 1.8463 = \$185
\item
  Much better than the \$0 return under no coordination!
\end{itemize}

2. (Diversification of Risk through Financial Intermediation) In an
economy, there are many identically distributed and independent projects
with 1 year of the investment horizon, and the (net) rate of return on
each project follows a normal distribution, with mean of 6\% and
standard deviation of 15\%. Each project requires an initial investment
of \$1000. (For questions a)-c), you need to use the Normal Distribution
Table in the appendix or you can use Excel or other statistical
software.)

Each project: return \(r \sim N(\mu,\sigma^{2})\)where \(\mu = 6\%\),
\(\sigma = 15\%\), initial investment = \$1,000.

\hypertarget{part-a-single-project}{%
\subsection{Part (a): Single Project}\label{part-a-single-project}}

We need probabilities for losses exceeding \$500 and \$100.

\textbf{Lose more than \$500} means return \(\left. <-50\% \right.\ \):

\[z = \frac{- 0.50 - 0.06}{0.15} = \frac{- 0.56}{0.15} = - 3.733\]

\[P(r < - 50\%) = \Phi( - 3.733) \approx \mathbf{0.0094\%}\]

\textbf{Lose more than \$100} means return \(\left. <-10\% \right.\ \):

\[z = \frac{- 0.10 - 0.06}{0.15} = \frac{- 0.16}{0.15} = - 1.067\]

\[P(r < - 10\%) = \Phi( - 1.067) \approx \mathbf{14.31\%}\]

So there is a non-trivial 14.3\% chance a single-project bank loses more
than 10\% of its
investment.{[}\href{https://ppl-ai-file-upload.s3.amazonaws.com/web/direct-files/attachments/64012734/8ecdf432-f8fa-409d-b72d-017134eed977/Homework1.pdf}{ppl-ai-file-upload.s3.amazonaws}{]}\hspace{0pt}

\hypertarget{part-b-portfolio-of-15-projects}{%
\subsection{Part (b): Portfolio of 15
Projects}\label{part-b-portfolio-of-15-projects}}

By the Central Limit Theorem, the average return of 15 i.i.d. projects
is:

\[\ \overline{}r_{15} \sim N\left( \mu,\frac{\sigma}{N} \right)\]

\begin{itemize}
\item
  \textbf{Mean} = \textbf{6\%} (unchanged)
\item
  \textbf{Standard Deviation} =
  \(\frac{15\%}{\sqrt{15}} = \mathbf{3.873\%}\)
\end{itemize}

\textbf{Lose more than \$1,500} on a \$15,000 total investment means
average return
\(\left. <-10\% \right.\ \):{[}\href{https://ppl-ai-file-upload.s3.amazonaws.com/web/direct-files/attachments/64012734/8ecdf432-f8fa-409d-b72d-017134eed977/Homework1.pdf}{ppl-ai-file-upload.s3.amazonaws}{]}\hspace{0pt}

\[z = \frac{- 0.10 - 0.06}{0.03873} = - 4.131\]

\[P({\overset{ˉ}{r}}_{15} < - 10\%) = \Phi( - 4.131) \approx \mathbf{0.00180\%}\]

Diversifying across 15 projects reduces the probability of a 10\% loss
from \textbf{14.3\%} to just \textbf{0.0018\%} --- a dramatic reduction
in
risk.{[}\href{https://ppl-ai-file-upload.s3.amazonaws.com/web/direct-files/attachments/64012734/8ecdf432-f8fa-409d-b72d-017134eed977/Homework1.pdf}{ppl-ai-file-upload.s3.amazonaws}{]}\hspace{0pt}

\hypertarget{part-c-n-projects-general-case}{%
\subsection{Part (c): N Projects (General
Case)}\label{part-c-n-projects-general-case}}

For a portfolio of \(N\)independent, identically distributed
projects:{[}\href{https://ppl-ai-file-upload.s3.amazonaws.com/web/direct-files/attachments/64012734/8ecdf432-f8fa-409d-b72d-017134eed977/Homework1.pdf}{ppl-ai-file-upload.s3.amazonaws}{]}\hspace{0pt}

\[{\overset{ˉ}{r}}_{N} \sim N\text{ ⁣}\left( \mu,\text{\ }\text{	}\text{	}\text{	}\text{	}\text{	}\text{	}\text{	}\frac{\sigma^{2}}{N} \right)\]

\begin{longtable}[]{@{}
  >{\raggedright\arraybackslash}p{(\columnwidth - 4\tabcolsep) * \real{0.2571}}
  >{\raggedright\arraybackslash}p{(\columnwidth - 4\tabcolsep) * \real{0.1646}}
  >{\raggedright\arraybackslash}p{(\columnwidth - 4\tabcolsep) * \real{0.5783}}@{}}
\toprule()
\begin{minipage}[b]{\linewidth}\raggedright
\textbf{Statistic}
\end{minipage} & \begin{minipage}[b]{\linewidth}\raggedright
\textbf{Formula}
\end{minipage} & \begin{minipage}[b]{\linewidth}\raggedright
\textbf{Value}
\end{minipage} \\
\midrule()
\endhead
Mean & \(\mu\) & \textbf{6\%} (constant, independent of N) \\
Std. Deviation & \(\sigma/\sqrt{N}\) & \(15\%/\sqrt{N}\)→ \textbf{0} as
\(N \rightarrow \infty\) \\
\bottomrule()
\end{longtable}

\textbf{As} \(N \rightarrow \infty\)\textbf{:} The standard deviation
collapses to zero, and by the Law of Large Numbers, the portfolio return
converges \textbf{almost surely} to the mean:

\[{\overset{ˉ}{r}}_{N}\overset{\phantom{a.s.}}{\rightarrow}\mu = 6\%\]

This is the core insight of financial intermediation: a bank holding
infinitely many uncorrelated projects faces \textbf{zero variance risk}
--- the return becomes a certainty of 6\%.

\hypertarget{part-d-maximum-cd-rate-on-a-competitive-bank}{%
\subsection{Part (d): Maximum CD Rate on a Competitive
Bank}\label{part-d-maximum-cd-rate-on-a-competitive-bank}}

When \(N \rightarrow \infty\), the bank earns exactly \(\mu = 6\%\)on
its asset portfolio with \textbf{no
uncertainty}.{[}\href{https://ppl-ai-file-upload.s3.amazonaws.com/web/direct-files/attachments/64012734/8ecdf432-f8fa-409d-b72d-017134eed977/Homework1.pdf}{ppl-ai-file-upload.s3.amazonaws}{]}\hspace{0pt}

A \textbf{fully competitive bank} earns zero profit, so it passes the
entire return to depositors:

\[\text{Maximum\ CD\ rate} = \mu = \boxed{6\%}\]

\textbf{Intuition:} The bank has transformed risky individual loans
(each with 15\% standard deviation) into a riskless liability --- a
fixed-return CD --- by pooling and diversifying. This perfectly
illustrates how banks create \textbf{safe secondary securities}
(deposits) from \textbf{risky primary assets} (loans), which is a
fundamental function of financial
intermediation.{[}\href{https://ppl-ai-file-upload.s3.amazonaws.com/web/direct-files/attachments/64012734/8ecdf432-f8fa-409d-b72d-017134eed977/Homework1.pdf}{ppl-ai-file-upload.s3.amazonaws}{]}\hspace{0pt}

3. (Bank Run) Consider a simplified version of the model developed in
Diamond and Dybvig (1983).

A bank collects deposits of 100 from many identical depositors at date
0.

At date 1, a fraction 20\% of depositors are impatient and must consume
early.

The remaining 80\% are patient and can wait until date 2.

The bank invests all funds in a long-term project that yields:

1.4 per unit if held to date 2

0.8 per unit if liquidated early at date 1

The bank offers a demand deposit contract promising:

1.0 if withdrawn at date 1

1.2 if withdrawn at date 2

\hypertarget{solution-to-question-3-bank-run-diamond-dybvig-model}{%
\subsection{Solution to Question 3: Bank Run (Diamond-Dybvig
Model)}\label{solution-to-question-3-bank-run-diamond-dybvig-model}}

\hypertarget{setup-summary}{%
\subsection{Setup Summary}\label{setup-summary}}

\begin{longtable}[]{@{}
  >{\raggedright\arraybackslash}p{(\columnwidth - 2\tabcolsep) * \real{0.6501}}
  >{\raggedright\arraybackslash}p{(\columnwidth - 2\tabcolsep) * \real{0.3499}}@{}}
\toprule()
\begin{minipage}[b]{\linewidth}\raggedright
\textbf{Item}
\end{minipage} & \begin{minipage}[b]{\linewidth}\raggedright
\textbf{Value}
\end{minipage} \\
\midrule()
\endhead
Total deposits collected (date 0) & 100 \\
Impatient depositors (withdraw date 1) & 20\% → 20 depositors \\
Patient depositors (withdraw date 2) & 80\% → 80 depositors \\
Long-term yield (held to date 2) & 1.4 per unit \\
Early liquidation yield (date 1) & 0.8 per unit \\
Demand deposit promise (date 1) & 1.0 per depositor \\
Demand deposit promise (date 2) & 1.2 per depositor \\
\bottomrule()
\end{longtable}

\hypertarget{part-a-solvency-vs.-liquidity}{%
\subsection{Part (a): Solvency vs.
Liquidity}\label{part-a-solvency-vs.-liquidity}}

\textbf{Scenario:} Only the 20 impatient depositors withdraw at date 1.

\textbf{Date 1 obligations:}

\begin{itemize}
\item
  20 depositors × 1.0 = \textbf{20 units} needed
\end{itemize}

\textbf{To meet this, the bank liquidates early:}

\begin{itemize}
\item
  Assets needed to liquidate = 20 / 0.8 = \textbf{25 units of
  investment}
\item
  Remaining investment = 100 − 25 = \textbf{75 units}
\end{itemize}

\textbf{Date 2 obligations:}

\begin{itemize}
\item
  80 patient depositors × 1.2 = \textbf{96 units} needed
\end{itemize}

\textbf{Date 2 assets available:}

\begin{itemize}
\item
  75 units × 1.4 = \textbf{105 units}
\end{itemize}

\textbf{Conclusion: The bank is fully
solvent.}{[}\href{https://ppl-ai-file-upload.s3.amazonaws.com/web/direct-files/attachments/64012734/8ecdf432-f8fa-409d-b72d-017134eed977/Homework1.pdf}{ppl-ai-file-upload.s3.amazonaws}{]}\hspace{0pt}

\[\text{Date\ 2\ Assets} = 105 \geq 96 = \text{Date\ 2\ Obligations}\]

The bank meets all obligations and retains a surplus of \textbf{9
units}. This demonstrates a healthy bank is only a \emph{liquidity}
mismatch (short-term vs. long-term), not an \emph{insolvency} problem.

\hypertarget{part-b-bank-run-equilibrium}{%
\subsection{Part (b): Bank Run
Equilibrium}\label{part-b-bank-run-equilibrium}}

\hypertarget{i.-what-happens-if-patient-depositors-panic-and-withdraw-at-date-1}{%
\subsection{i. What happens if patient depositors panic and withdraw at
date
1?}\label{i.-what-happens-if-patient-depositors-panic-and-withdraw-at-date-1}}

If patient depositors fear others will run, their dominant strategy
becomes withdrawing early too. The bank now faces \textbf{all 100
depositors} demanding 1.0 at date 1, but its assets yield only
\textbf{0.8 per unit} when liquidated early. The bank cannot honour its
promise of 1.0 per depositor and will \textbf{fail} before all
depositors are
paid.{[}\href{https://ppl-ai-file-upload.s3.amazonaws.com/web/direct-files/attachments/64012734/8ecdf432-f8fa-409d-b72d-017134eed977/Homework1.pdf}{ppl-ai-file-upload.s3.amazonaws}{]}\hspace{0pt}

\hypertarget{ii.-can-the-bank-survive-if-80-all-patient-depositors-withdraw-at-date-1}{%
\subsection{ii. Can the bank survive if 80\% (all patient depositors)
withdraw at date
1?}\label{ii.-can-the-bank-survive-if-80-all-patient-depositors-withdraw-at-date-1}}

Total withdrawals demanded at date 1:

\begin{itemize}
\item
  All 100 depositors × 1.0 = \textbf{100 units} needed
\end{itemize}

Total assets available if liquidated early:

\begin{itemize}
\item
  100 units × 0.8 = \textbf{80 units}
\end{itemize}

\textbf{Shortfall = 80 − 100 = −20 units}

The bank can only pay \textbf{80 cents per dollar} promised. Using a
sequential service constraint (first-come,
first-served):{[}\href{https://ppl-ai-file-upload.s3.amazonaws.com/web/direct-files/attachments/64012734/8ecdf432-f8fa-409d-b72d-017134eed977/Homework1.pdf}{ppl-ai-file-upload.s3.amazonaws}{]}\hspace{0pt}

\begin{itemize}
\item
  First \textbf{80 depositors} in line receive \textbf{1.0 each} (total
  paid = 80)
\item
  Remaining \textbf{20 depositors} receive \textbf{\$0} (bank is
  insolvent)
\end{itemize}

\[\text{Assets} = 80 < 100 = \text{Obligations} \Rightarrow \text{Bank\ }\text{	}\text{	}\text{	}\text{	}\text{	fails}\]

\hypertarget{iii.-why-is-this-a-self-fulfilling-equilibrium}{%
\subsection{iii. Why is this a self-fulfilling
equilibrium?}\label{iii.-why-is-this-a-self-fulfilling-equilibrium}}

This is a \textbf{self-fulfilling prophecy} because the belief itself
causes the outcome it
predicts:{[}\href{https://ppl-ai-file-upload.s3.amazonaws.com/web/direct-files/attachments/64012734/8ecdf432-f8fa-409d-b72d-017134eed977/Homework1.pdf}{ppl-ai-file-upload.s3.amazonaws}{]}\hspace{0pt}

\begin{itemize}
\item
  \textbf{If} a patient depositor believes others will run → rational to
  withdraw early (get 1.0 now vs. risk getting 0 later)
\item
  \textbf{Because} enough people run → bank truly fails
\item
  \textbf{Therefore} the fear was justified, validating the original
  belief
\end{itemize}

This creates \textbf{two Nash Equilibria}:

\begin{enumerate}
\def\labelenumi{\arabic{enumi}.}
\item
  \textbf{Good equilibrium}: Patient depositors wait → bank is solvent →
  everyone is paid (1.0 at date 1, 1.2 at date 2)
\item
  \textbf{Bad equilibrium (bank run)}: Everyone withdraws at date 1 →
  bank fails → self-confirming
\end{enumerate}

Neither equilibrium depends on the bank\textquotesingle s
\emph{fundamental} health --- the bad equilibrium is driven purely by
\textbf{coordination failure and panic}, not insolvency. This is Diamond
and Dybvig\textquotesingle s central insight.

\hypertarget{part-c-policy-interventions}{%
\subsection{Part (c): Policy
Interventions}\label{part-c-policy-interventions}}

\hypertarget{i.-how-does-deposit-insurance-change-the-equilibrium}{%
\subsection{i. How does deposit insurance change the
equilibrium?}\label{i.-how-does-deposit-insurance-change-the-equilibrium}}

Deposit insurance eliminates the bad equilibrium entirely by removing
the incentive to run. With government-guaranteed
deposits:{[}\href{https://ppl-ai-file-upload.s3.amazonaws.com/web/direct-files/attachments/64012734/8ecdf432-f8fa-409d-b72d-017134eed977/Homework1.pdf}{ppl-ai-file-upload.s3.amazonaws}{]}\hspace{0pt}

\begin{itemize}
\item
  A patient depositor knows they will receive \textbf{1.2 at date 2
  regardless} of what others do
\item
  Withdrawing early yields only 1.0 \textless{} 1.2, so waiting is
  always dominant
\item
  \textbf{No rational depositor runs} → the bad equilibrium ceases to
  exist
\end{itemize}

The key mechanism is that deposit insurance \textbf{breaks the strategic
complementarity}: your optimal action (wait) no longer depends on what
you fear others will do. The cost is moral hazard --- banks may take
excessive risks knowing deposits are guaranteed --- which is why deposit
insurance is typically paired with prudential regulation and capital
requirements.

\hypertarget{ii.-would-a-lender-of-last-resort-lolr-eliminate-the-bad-equilibrium}{%
\subsection{ii. Would a lender of last resort (LOLR) eliminate the bad
equilibrium?}\label{ii.-would-a-lender-of-last-resort-lolr-eliminate-the-bad-equilibrium}}

A central bank acting as \textbf{lender of last resort}
(Bagehot\textquotesingle s doctrine) can also eliminate the bad
equilibrium, \textbf{but only under specific
conditions}:{[}\href{https://ppl-ai-file-upload.s3.amazonaws.com/web/direct-files/attachments/64012734/8ecdf432-f8fa-409d-b72d-017134eed977/Homework1.pdf}{ppl-ai-file-upload.s3.amazonaws}{]}\hspace{0pt}

\begin{longtable}[]{@{}
  >{\raggedright\arraybackslash}p{(\columnwidth - 2\tabcolsep) * \real{0.3730}}
  >{\raggedright\arraybackslash}p{(\columnwidth - 2\tabcolsep) * \real{0.6270}}@{}}
\toprule()
\begin{minipage}[b]{\linewidth}\raggedright
\textbf{Condition}
\end{minipage} & \begin{minipage}[b]{\linewidth}\raggedright
\textbf{Why it matters}
\end{minipage} \\
\midrule()
\endhead
Lend freely at a penalty rate & Ensures only solvent-but-illiquid banks
borrow, not insolvent ones \\
Against good collateral & Limits central bank risk; forces banks to hold
quality assets \\
The bank must be \textbf{illiquid, not insolvent} & LOLR cures a
\emph{liquidity} crisis, not a \emph{solvency} crisis \\
Credibly communicated & The mere \textbf{announcement} of LOLR support
can prevent a run from starting \\
\bottomrule()
\end{longtable}

In this model\textquotesingle s context, the bank \emph{is} solvent (as
shown in Part a), so a LOLR pledge that "the central bank will provide
liquidity if needed" would make it irrational for patient depositors to
run --- they know the bank can always honour 1.2 at date 2. If the bank
were truly insolvent, however, LOLR support would merely delay failure
and create moral hazard without resolving the underlying problem.

\emph{Note: This solution was developed with AI assistance for numerical
verification, with analysis grounded in Diamond and Dybvig (1983)
framework.}

\end{document}
